\documentclass[../../../include/open-logic-section]{subfiles}

\begin{document}

\olfileid[id]{sth}{card-arithmetic}{opps}
\olsection{Mendefinisikan Operasi Dasar}

Karena kita tidak perlu melacak urutan, aritmetika kardinal agak lebih mudah
didefinisikan daripada aritmetika ordinal. Kita akan mendefinisikan
penjumlahan, perkalian, dan eksponensiasi secara bersamaan.

\begin{defn}
Jika $\cardfont{a}$ dan $\cardfont{b}$ merupakan kardinal:
\begin{align*}
	\cardfont{a} \cardplus \cardfont{b} &\defis 
	\card{\cardfont{a} \disjointsum \cardfont{b}}\\
	\cardfont{a} \cardtimes \cardfont{b} &\defis 
	\card{\cardfont{a} \times \cardfont{b}}\\
	\cardexpo{\cardfont{a}}{\cardfont{b}} &\defis 
	\card{\funfromto{\cardfont{b}}{\cardfont{a}}}
\end{align*}
di mana $\funfromto{X}{Y} = \Setabs{f}{f \text{ merupakan fungsi } X \to Y}$.
(Mudah ditunjukkan bahwa $\funfromto{X}{Y}$ ada untuk sebarang himpunan $X$
dan $Y$; kita menyerahkannya sebagai latihan.)
\end{defn}

\begin{prob}
Buktikan dalam $\Zminus$ bahwa $\funfromto{X}{Y}$ ada untuk sebarang himpunan
$X$ dan~$Y$. Dengan bekerja dalam $\ZF$, hitung
$\setrank{\funfromto{X}{Y}}$ dari $\setrank{X}$ dan $\setrank{Y}$ dengan cara
seperti dalam \olref[sth][ord-arithmetic][using-addition]{rankcomputation}.
\end{prob}

Penjelasan berikut mungkin membantu memahami definisi ini. Mengenai
penjumlahan: definisi tersebut menggunakan gagasan jumlah lepas,
$\disjointsum$, sebagaimana didefinisikan dalam
\olref[ord-arithmetic][add]{defdissum}; mudah dilihat bahwa definisi ini
memberikan hasil yang benar untuk kasus-kasus hingga. Mengenai perkalian:
\olref[sfr][set][pai]{cardnmprod} memberi tahu kita bahwa jika $A$ mempunyai
$n$ anggota dan $B$ mempunyai $m$ anggota, maka $A \times B$ mempunyai
$n \cdot m$ anggota; jadi, definisi kita sekadar menggeneralisasi gagasan
tersebut menjadi perkalian transfinit. Eksponensiasi serupa: kita hanya
menggeneralisasi gagasan dari kasus hingga ke kasus transfinit. Bahkan, dalam
beberapa hal, aritmetika kardinal transfinit jauh lebih menyerupai aritmetika
``biasa'' daripada aritmetika ordinal transfinit:

\begin{prop}\ollabel{cardplustimescommute}
$\cardplus$ dan $\cardtimes$ bersifat komutatif dan asosiatif.
\end{prop}

\begin{proof}
Untuk komutativitas, berdasarkan
\olref[cardinals][cardsasords]{lem:CardinalsBehaveRight}, cukup diamati bahwa
$\cardeq{(\cardfont{a} \disjointsum
\cardfont{b})}{(\cardfont{b} \disjointsum \cardfont{a})}$ dan
$\cardeq{(\cardfont{a} \times \cardfont{b})}{(\cardfont{b} \times
\cardfont{a})}$. Kita menyerahkan asosiativitas sebagai latihan.
\end{proof}

\begin{prob}
Buktikan bahwa $\cardplus$ dan $\cardtimes$ bersifat asosiatif.
\end{prob}

\begin{prop}
$A$ bersifat tak hingga jika dan hanya jika
$\card{A} \cardplus 1 = 1 \cardplus \card{A} = \card{A}$.
\end{prop}

\begin{proof}
Seperti dalam \olref[cardinals][classing]{generalinfinitycharacter}, hal ini
mengikuti dari
\olref[ord-arithmetic][using-addition]{ordinfinitycharacter} dan
\olref[cardinals][cardsasords]{lem:CardinalsBehaveRight}.
\end{proof}

Hal ini menjelaskan mengapa kita perlu menggunakan simbol yang berbeda bagi
penjumlahan/perkalian ordinal dan kardinal: keduanya sungguh-sungguh
merupakan operasi yang \emph{berbeda}. Pasangan hasil berikut menunjukkan
bahwa eksponensiasi ordinal dan kardinal juga merupakan operasi yang berbeda.
(Ingat bahwa \olref[z][infinity-again]{defnomega} menyiratkan bahwa
$2 = \{0, 1\}$):

\begin{lem}\ollabel{lem:SizePowerset2Exp}
$\card{\Pow{A}} = \cardexpo{2}{\card{A}}$, untuk sebarang $A$.
\end{lem}

\begin{proof}
Untuk setiap himpunan bagian $B \subseteq A$, misalkan
$\chi_B \in \funfromto{A}{2}$ diberikan oleh:
\begin{align*}
	\chi_{B}(x) &\defis
	\begin{cases}
		1 & \text{jika }x\in B\\
		0 & \text{jika tidak.}
	\end{cases}
\end{align*}
Sekarang misalkan $f(B) = \chi_B$; ini mendefinisikan !!a{bijection}
$f \colon \Pow{A} \to \funfromto{A}{2}$. Jadi
$\cardeq{\Pow{A}}{\funfromto{A}{2}}$. Dengan demikian,
$\cardeq{\Pow{A}}{\funfromto{\card{A}}{2}}$, sehingga
$\card{\Pow{A}} = \card{\funfromto{\card{A}}{2}} =
2^{\card{A}}$.
\end{proof}

Bukti ringkas ini pada dasarnya mencakup pembahasan dalam
\olref[sfr][siz][red-alt]{sec}. Di sana, kita menunjukkan cara
``mereduksi'' ketakterhitungan $\Pow{\omega}$ menjadi ketakterhitungan
himpunan semua string biner tak hingga, $\Bin^\omega$. Pada dasarnya,
$\Bin^{\omega}$ hanyalah $\funfromto{\omega}{2}$; dan bukti sebelumnya
menunjukkan bahwa penalaran yang kita gunakan dalam
\olref[sfr][siz][red-alt]{sec} akan tetap berlaku dengan menggunakan
sebarang himpunan~$A$ sebagai pengganti~$\omega$. Hasil tersebut juga
memberikan suatu fakta singkat tentang eksponensiasi kardinal:

\begin{cor}\ollabel{cantorcor}
$\cardfont{a} < \cardexpo{2}{\cardfont{a}}$ untuk sebarang
kardinal~$\cardfont{a}$.
\end{cor}

\begin{proof}
Dari Teorema Cantor (\olref[sfr][siz][car]{thm:cantor}) dan
\olref{lem:SizePowerset2Exp}.
\end{proof}
\noindent
Jadi $\omega < \cardexpo{2}{\omega}$. Namun, perhatikan bahwa ini merupakan
hasil tentang eksponensiasi \emph{kardinal}. Hasil ini perlu dibandingkan
dengan eksponensiasi \emph{ordinal}, sebab dalam kasus kedua berlaku
$\omega = \ordexpo{2}{\omega}$ (lihat
\olref[ord-arithmetic][expo]{sec}).

Selagi membahas eksponensiasi kardinal, kita juga dapat menjelaskan secara
sedikit lebih tepat ``cara'' $\Real$ bersifat !!{nonenumerable}.

\begin{thm}\ollabel{continuumis2aleph0}
$\card{\Real} = \cardexpo{2}{\omega}$
\end{thm}

\begin{proof}[Kerangka bukti]
Terdapat banyak cara untuk membuktikannya. Cara yang paling langsung ialah
menunjukkan bahwa $\cardle{\Pow{\omega}}{\Real}$ dan
$\cardle{\Real}{\Pow{\omega}}$, lalu menggunakan Schr\"oder--Bernstein untuk
menyimpulkan bahwa $\cardeq{\Real}{\Pow{\omega}}$, dan menggunakan
\olref[card-arithmetic][opps]{lem:SizePowerset2Exp} untuk menyimpulkan bahwa
$\card{\Real} = \cardexpo{2}{\omega}$. Kita menyerahkan sebagai latihan
(yang mencerahkan) untuk mendefinisikan injeksi
$f \colon \Pow{\omega} \to \Real$ dan
$g \colon \Real \to \Pow{\omega}$.
\end{proof}

\begin{prob}
Lengkapi bukti \olref[sth][card-arithmetic][opps]{continuumis2aleph0} dengan
menunjukkan bahwa $\cardle{\Pow{\omega}}{\Real}$ dan
$\cardle{\Real}{\Pow{\omega}}$.
\end{prob}

\end{document}
